\emph{Setting.} $e$ is orthogonal to every $e_{k}$, so $e$ anticommutes with each $e_{k}$ and $e_{1}, \dots, e_{n}, e$ is an orthonormal basis of $W$, of odd dimension $n + 1$. Every basis blade of $\Cl(W)$ either omits $e$, and is a blade of $\Cl(V)$, or contains it once, and is then a blade of $\Cl(V)$ times $e$; so $\Cl(W) = \Cl(V) \oplus \Cl(V)e$ as real vector spaces, of dimensions $2^{n}$ each.

\emph{(i).} Since $n$ is even, $e$ anticommutes with each of the $n$ factors of $I_{V}$ and so commutes with $I_{V}$, and every $e_{k}$ anticommutes with the $n - 1$ other factors and commutes with its own, so $e_{k}I_{V} = (-1)^{n-1}I_{V}e_{k} = -I_{V}e_{k}$. Hence $I_{W}^{2} = I_{V}eI_{V}e = I_{V}^{2}e^{2} = \varepsilon_{V}Q(e)$, which is $-\varepsilon_{V}^{2} = -1$ for $Q(e) = -\varepsilon_{V}$ and $+1$ for $Q(e) = +\varepsilon_{V}$. The value $\varepsilon_{V} = (-1)^{n(n-1)/2 + q}$ is the reordering sign of $e_{1}\cdots e_{n}e_{1}\cdots e_{n}$ times the product of the squares. $I_{W}$ is the volume element of $W$, of odd dimension, so it is central with square $(-1)^{(n+1)n/2 + q'}$, $q'$ the number of negative squares of $W$ (Theorem \claimref{odd-extensions.volume-element-odd-dimension}); directly, $e_{k}I_{W} = e_{k}I_{V}e = -I_{V}e_{k}e = I_{V}ee_{k} = I_{W}e_{k}$ and $eI_{W} = eI_{V}e = I_{V}e^{2} = I_{W}e$, so $I_{W}$ commutes with every vector of $W$ and, the vectors generating $\Cl(W)$, it is central. The centre of $\Cl(W)$, $\dim W$ odd, is $\R \cdot 1 \oplus \R \cdot I_{W}$ (Theorem \claimref{primer.centre}). A real algebra $\mathrm{span}\{1, u\}$ with $u^{2} = -1$ is a field isomorphic to $\C$: $(x + yu)(x - yu) = x^{2} + y^{2}$ is a nonzero real number unless $x = y = 0$. The script computes $\varepsilon_{V}$, $I_{W}^{2} = -1$ with $Q(e) = -\varepsilon_{V}$ and $+1$ with the other sign, and the commutation of $I_{W}$ with every generator, for every signature of dimension $2$ and $4$ and for $\Cl(2,4)$, and finds the commutant of the generators of dimension exactly $2$ for dimensions $2$ and $4$.

\emph{(ii).} $I_{V}$ is invertible ($I_{V}^{-1} = \varepsilon_{V}I_{V}$), so $\Cl(V)e = \Cl(V)I_{V}e = \Cl(V)I_{W}$, and with the setting $\Cl(W) = \Cl(V) \oplus \Cl(V)I_{W}$. The multiplication map $m: a \otimes (x + yI_{W}) \mapsto a(x + yI_{W})$ is therefore a linear bijection $\Cl(V) \otimes_{\R} \mathrm{span}\{1, I_{W}\} \to \Cl(W)$, and it is multiplicative because $I_{W}$ is central: $m(a \otimes z)\,m(b \otimes w) = azbw = abzw = m(ab \otimes zw)$. The complexification of $\Cl(V)$ is $\Cl(V) \otimes_{\R} F$ for a field $F \cong \C$ with its unit $u$, presented by the generators $e_{k}$ with the Clifford relations and one central $u$ with $u^{2} = -1$; the images $e_{k}$ and $I_{W}$ satisfy exactly these relations, so the map on generators extends to a homomorphism, which is $m$ and so an isomorphism. By Theorem \claimref{std-algebra.complexification}(i) the complexification of $\Cl(V)$, $n$ even, is a full matrix algebra of size $2^{n/2}$ over the complex field, here $\mathrm{span}\{1, I_{W}\}$. Since $\dim W$ is odd and $I_{W}^{2} = -1$, Theorem \claimref{odd-extensions.doubling-or-complexification}(ii) gives $\Cl(W) \cong \Cl^{0}(W) \otimes_{\R} \mathrm{span}\{1, I_{W}\}$; concretely $I_{W}$ is odd and invertible, so $\Cl^{0}(W)I_{W}$ is the odd part and $\Cl(W) = \Cl^{0}(W) \oplus \Cl^{0}(W)I_{W}$. The script checks, for $\Cl(1,3) \to \Cl(2,3)$ and $\Cl(2,4) \to \Cl(3,4)$, the relations of the images of the generators, that the images of the real basis $\{x, xI_{W}\}$ have rank $2^{n+1}$, the even-odd splitting, and, for $\Cl(1,3) \to \Cl(2,3)$, that $m$ is multiplicative on all $32 \times 32$ pairs of basis elements, the abstract product of pairs $(a, b) \leftrightarrow a + bu$ being $(a, b)(c, d) = (ac - bd, ad + bc)$.

\emph{(iii).} $I_{W}$ commutes with $a$ and $b$ and $I_{W}^{2} = -1$, so $(a + I_{W}b)^{2} = a^{2} + I_{W}(ab + ba) + I_{W}^{2}b^{2} = Q(a) - Q(b) + 2\ip{a}{b}I_{W}$, using $ab + ba = 2\ip{a}{b}$. This is the complex-bilinear extension of $Q$: it is $Q(a) - Q(b)$ plus $I_{W}$ times twice the polar form, as $Q(a + \iota b) = Q(a) + 2\iota\ip{a}{b} + \iota^{2}Q(b)$ is for any unit $\iota$ of square $-1$ commuting with $V$. It vanishes exactly when its two real coordinates do, $Q(a) = Q(b)$ and $\ip{a}{b} = 0$, since $1$ and $I_{W}$ are independent. With $a \to \epsilon a$ and $y$ in place of $a$: $(y + \epsilon I_{W}a)^{2} = Q(y) - \epsilon^{2}Q(a) + 2\epsilon\ip{y}{a}I_{W}$. The value lies in the commutative field $\mathrm{span}\{1, I_{W}\}$ by (i), so a polynomial, a quotient (by (i), every nonzero element of the field is invertible in it) or a convergent power series in that square is computed entirely inside the field, where $I_{W}$ is a central square root of $-1$ and every rule of complex arithmetic holds with $I_{W}$ as the unit. The script checks the square on three pairs of rational vectors in $\Cl(2,3)$, the null vector $e_{2} + I_{5}e_{3}$, the contrast $(e_{2} + I_{4}e_{3})^{2} = -2 + 2e_{0}e_{1}$ (the left $I_{4}$ anticommutes with vectors and gives no complexification), and the shift with its inverse in the field at two rational points.

\emph{(iv).} With $Q(e) = +1$, $P_{\pm}^{2} = \tfrac14(1 \pm 2e + e^{2}) = P_{\pm}$, and $e$ commutes with $I_{V}$ (as in (i)), so $P_{\pm}$ commutes with $I_{V}$. The map $x \mapsto xP_{\pm} = \tfrac12 x \pm \tfrac12 xe$ commutes with left multiplication by $\Cl(V)$, and it is injective because $x$ and $xe$ lie in the complementary summands $\Cl(V)$ and $\Cl(V)e$. Its image is $\Cl(V)P_{\pm}$, and since $eP_{\pm} = \pm P_{\pm}$, $\Cl(W)P_{\pm} = \Cl(V)P_{\pm} + \Cl(V)eP_{\pm} = \Cl(V)P_{\pm}$. Now $P_{\pm}I_{W} = P_{\pm}I_{V}e = I_{V}P_{\pm}e = \pm I_{V}P_{\pm}$, so $(xP_{\pm})I_{W} = \pm(xI_{V})P_{\pm}$. Since $n$ is even, $I_{V}$ commutes with even elements of $\Cl(V)$ and anticommutes with vectors, which gives the two further forms. The script checks the identity for both idempotents on all $16$ blades of $\Cl(1,3)$ and all $64$ blades of $\Cl(2,4)$, the even and vector forms, and that the images $xP_{+}$ of the blades have rank $2^{n}$.

\emph{(v).} $\Cl(1,3)$ has $n = 4$, $q = 3$, $\varepsilon_{V} = (-1)^{6 + 3} = -1$; $\Cl(2,4)$ with basis $e_{0}, \dots, e_{3}, e_{+}, e_{-}$ has $n = 6$, $q = 4$, $\varepsilon_{V} = (-1)^{15 + 4} = -1$, and its pseudoscalar in this order is $I_{4}e_{+}e_{-} = I_{6}$. In both $Q(e) = +1$ and the extension has one more positive square: $\Cl(2,3)$ and $\Cl(3,4)$. The null vector and the shift are (iii) with $Q(e_{2}) = Q(e_{3}) = -1$, $\ip{e_{2}}{e_{3}} = 0$, $a = e_{0}$, $Q(e_{0}) = 1$; the relation to $I_{4}$ and $I_{6}$ is (iv) with $P_{+}$.

\emph{(vi).} $V' = V \oplus \R f$ has even dimension, so (i) and (ii) apply to it: $I_{W'}$ is central in $\Cl(W')$ with square $-1$, and $m: a \otimes z \mapsto az$ restricted to $\Cl(V) \otimes_{\R} \mathrm{span}\{1, I_{W'}\}$ is multiplicative by centrality. It is injective: $I_{W'}$ is a top blade containing $f$ and $e$, so every blade of $\Cl(V)I_{W'}$ contains $f$ and $e$ while no blade of $\Cl(V)$ does, and $a + bI_{W'} = 0$ forces $a = 0$ and $bI_{W'} = 0$, hence $b = 0$. The script checks, for $V = \R^{1,2}$ with $Q(f) = -1$ and $Q(e) = +1$ (so $W' = \R^{2,3}$), multiplicativity on all $16 \times 16$ pairs of basis elements and rank $16$. $\square$
